\documentclass[11pt]{article}
\usepackage[margin=1in]{geometry}
\usepackage{enumitem}
% =============================================================================
% Preambles/header.tex — shared LaTeX/Beamer preamble for this project
%
% Use from a Beamer lecture by prepending:
%     \input{header}
% and compiling with TEXINPUTS=../Preambles:$TEXINPUTS (the /compile-latex
% skill does this automatically).
%
% PALETTE CONTRACT. Color names defined here MUST match the names in
% Quarto/theme-template.scss so Beamer and Quarto renderings use the same
% palette. The scripts/check-palette-sync.sh script enforces this.
%
% To customize: edit the HEX values below AND the matching $variable in
% Quarto/theme-template.scss, then run scripts/check-palette-sync.sh.
% =============================================================================

% -----------------------------------------------------------------------------
% Engine detection + encoding
%
% Our /compile-latex skill uses XeLaTeX, where inputenc is unnecessary and
% can produce encoding warnings. Only load inputenc under pdfTeX.
% -----------------------------------------------------------------------------
\usepackage{iftex}
\ifPDFTeX
  \usepackage[utf8]{inputenc}
\fi

\usepackage{amsmath, amssymb, amsthm}
\usepackage{booktabs}
\usepackage{graphicx}

% -----------------------------------------------------------------------------
% PALETTE — mirrors Quarto/theme-template.scss variable names.
% Load xcolor and define colors BEFORE anything that references them
% (notably hyperref, hypersetup, and the Beamer color assignments).
% -----------------------------------------------------------------------------
\usepackage{xcolor}

\definecolor{primary-blue}{HTML}{012169}      % headings, accents
\definecolor{primary-gold}{HTML}{B9975B}      % emphasis, borders
\definecolor{highlight-yellow}{HTML}{F2A900}  % markers, alerts
\definecolor{light-bg}{HTML}{E8EDF5}          % title / section backgrounds
\definecolor{jet}{HTML}{1A1A1A}               % body text

% Semantic colors (used in diagrams and annotations)
\definecolor{positive}{HTML}{15803D}          % good / observed / identified
\definecolor{negative}{HTML}{B91C1C}          % bad / problematic / bias
\definecolor{neutral}{HTML}{525252}           % reference / context

% Highlight accents (match .hi-* classes in the SCSS)
\definecolor{hi-slate}{HTML}{314F4F}
\definecolor{hi-green}{HTML}{2E8B57}
\definecolor{hi-red}{HTML}{C41E3A}

% Semantic aliases so skill templates and snippets can write
% \textcolor{accent}{...} without hard-coding "primary-blue".
\colorlet{accent}{primary-blue}
\colorlet{accent2}{primary-gold}

% -----------------------------------------------------------------------------
% Hyperref — loaded AFTER xcolor + palette so link colors resolve.
% -----------------------------------------------------------------------------
\usepackage{hyperref}
\hypersetup{colorlinks=true, linkcolor=primary-blue, urlcolor=primary-blue,
            citecolor=primary-blue, pdfborder={0 0 0}}

% -----------------------------------------------------------------------------
% Course code — set once per deck (after \input{header}, before \begin{document})
% via \coursecode{CS401}. Shown in the footer on every slide so a deck's course
% is identifiable without opening the title page. Empty by default.
% -----------------------------------------------------------------------------
\makeatletter
\newcommand{\coursecode}[1]{\gdef\@coursecodetext{#1}}
\gdef\@coursecodetext{}
\makeatother

% -----------------------------------------------------------------------------
% Beamer theme assignments (only apply under Beamer).
%
% We detect Beamer via \@ifclassloaded{beamer}, which is the documented,
% non-internal way. This must be wrapped in \makeatletter/\makeatother
% because \@ifclassloaded contains an @-catcode character.
% -----------------------------------------------------------------------------
\makeatletter
\@ifclassloaded{beamer}{%
  \usefonttheme{professionalfonts}
  \setbeamercolor{structure}{fg=primary-blue}
  \setbeamercolor{frametitle}{fg=primary-blue}
  \setbeamercolor{title}{fg=primary-blue}
  \setbeamercolor{subtitle}{fg=primary-gold}
  \setbeamercolor{author}{fg=primary-blue}
  \setbeamercolor{institute}{fg=neutral}
  \setbeamercolor{date}{fg=primary-gold}
  \setbeamercolor{itemize item}{fg=primary-blue}
  \setbeamercolor{itemize subitem}{fg=primary-gold}
  \setbeamercolor{alerted text}{fg=primary-gold}
  \setbeamercolor{block title}{fg=white, bg=primary-blue}
  \setbeamercolor{block body}{bg=light-bg}
  % Semantic block variants: block=definition/concept (blue), exampleblock=
  % worked example (green/positive), alertblock=key takeaway or caution (gold).
  % Keep to <=2 colored boxes per slide across ALL three variants (INV-7).
  \setbeamercolor{block title example}{fg=white, bg=positive}
  \setbeamercolor{block body example}{bg=light-bg}
  \setbeamercolor{block title alerted}{fg=white, bg=primary-gold}
  \setbeamercolor{block body alerted}{bg=light-bg}

  % Minimal footer: course code (left) + page X / N (right), no navigation clutter
  \setbeamertemplate{navigation symbols}{}
  \setbeamertemplate{footline}{%
    \color{neutral}\footnotesize\hspace{0.7em}\@coursecodetext\hfill
    \insertframenumber/\inserttotalframenumber\hspace{0.7em}\vspace{0.3em}}
}{}
\makeatother

% -----------------------------------------------------------------------------
% TikZ setup — libraries the snippet gallery uses
% -----------------------------------------------------------------------------
\usepackage{tikz}
\usetikzlibrary{arrows.meta, positioning, calc, decorations.pathreplacing,
                fit, shapes.geometric, backgrounds}

% Shared TikZ styles — snippets and hand-written diagrams can pick these up
% via `\begin{tikzpicture}[every node/.style={...}]` or by naming specific
% styles (e.g., `\node[dag-node] (X) at (0,0) {X};`).
\tikzset{
  % Node styles for causal DAGs and similar
  dag-node/.style = {
    draw=primary-blue, fill=light-bg, rounded corners=2pt,
    minimum width=1.2cm, minimum height=0.9cm, align=center, font=\small
  },
  decision-node/.style = {
    draw=primary-gold, fill=white, diamond, aspect=1.8,
    minimum width=1.8cm, minimum height=1.2cm, align=center, font=\small,
    inner sep=1pt
  },
  % Edge styles — observed vs counterfactual
  observed-edge/.style = {-{Latex[length=2.5mm]}, thick, draw=jet},
  counterfactual-edge/.style = {-{Latex[length=2.5mm]}, thick, draw=neutral, dashed},
  confound-edge/.style = {-{Latex[length=2.5mm]}, thick, draw=negative, dashed},
  % Data-point styles for plots
  observed-dot/.style = {fill=primary-blue, draw=primary-blue, circle, inner sep=1.2pt},
  counterfactual-dot/.style = {fill=white, draw=neutral, circle, inner sep=1.2pt}
}

% -----------------------------------------------------------------------------
% Convenience macros
% -----------------------------------------------------------------------------
\newcommand{\muted}[1]{\textcolor{neutral}{#1}}
\newcommand{\key}[1]{\textcolor{primary-gold}{\textbf{#1}}}
\newcommand{\good}[1]{\textcolor{positive}{#1}}
\newcommand{\bad}[1]{\textcolor{negative}{#1}}

% Standout frame macro: full-bleed dark background for section transitions.
% Beamer-only (uses \begin{frame}/\setbeamercolor/\usebeamerfont) -- do not
% call from an article-class file (e.g. Notes/<CODE>/*-notes.tex). \muted,
% \key, \good, \bad, and \coursecode above are plain xcolor/gdef and are
% safe to use from either class; verified when Notes/article-class support
% was added.
% Expand: \transitionslide{Your transition title}
\newcommand{\transitionslide}[1]{%
  \begingroup
  \setbeamercolor{background canvas}{bg=primary-blue}
  \begin{frame}[plain]
    \centering
    \vfill
    {\usebeamerfont{title}\color{white}\Large #1\par}
    \vfill
  \end{frame}
  \endgroup
}

% Section divider frame (Beamer-only), an alternative to \transitionslide with a
% darker two-line style: near-black (jet) background, a small white label line
% above a larger gold title, and a thin gold rule along the bottom edge.
% Expand: \sectiondivider{Section 2}{Pointers}
\newcommand{\sectiondivider}[2]{%
  \begingroup
  \setbeamercolor{background canvas}{bg=jet}
  \begin{frame}[plain]
    \vspace*{3.0cm}
    {\color{white}\ttfamily\large #1\par}
    \vspace{0.5cm}
    {\color{primary-gold}\LARGE #2\par}
    \begin{tikzpicture}[remember picture, overlay]
      \draw[primary-gold, line width=2.5pt]
        ($(current page.south west)+(0,0.1)$) -- ($(current page.south east)+(0,0.1)$);
    \end{tikzpicture}
  \end{frame}
  \endgroup
}

% -----------------------------------------------------------------------------
% Channel watermark — "Concepts That Click" (YouTube).
%
% Article-class files (CompetitiveExam Q/A sets, Notes, Assignments) get a
% light diagonal draft watermark on every page plus a centered footer naming
% the channel. Beamer decks get a subtle TikZ background watermark instead
% (the footline already carries the course code). Centralized here so every
% MCQ PDF is branded without per-file edits.
% -----------------------------------------------------------------------------
\makeatletter
\@ifclassloaded{article}{%
  \usepackage{draftwatermark}
  \SetWatermarkText{Concepts That Click}
  \SetWatermarkScale{1}
  \SetWatermarkFontSize{2cm}
  \SetWatermarkLightness{0.86}
  \SetWatermarkAngle{45}
  \usepackage{fancyhdr}
  \pagestyle{fancy}
  \fancyhf{}
  \fancyfoot[C]{\footnotesize\textcolor{neutral}{Concepts That Click~|~YouTube: \href{https://www.youtube.com/@concepts.thatclick}{@concepts.thatclick}}}
  \renewcommand{\headrulewidth}{0pt}
  \renewcommand{\footrulewidth}{0pt}
  \fancypagestyle{plain}{%
    \fancyhf{}%
    \fancyfoot[C]{\footnotesize\textcolor{neutral}{Concepts That Click~|~YouTube: \href{https://www.youtube.com/@concepts.thatclick}{@concepts.thatclick}}}%
    \renewcommand{\headrulewidth}{0pt}%
    \renewcommand{\footrulewidth}{0pt}%
  }
}{}
\@ifclassloaded{beamer}{%
  \setbeamertemplate{background}{%
    \begin{tikzpicture}[remember picture, overlay]
      \node[rotate=45, opacity=0.055, align=center]
        at (current page.center)
        {\Huge\textcolor{neutral}{Concepts That Click}};
    \end{tikzpicture}%
  }
}{}
\makeatother 
\coursecode{JKSSB}
\renewcommand{\thesection}{41.\arabic{section}}
\newtheorem{example}{Example}[section]
\newenvironment{definitionbox}[1]{%
  \par\noindent\textbf{Definition (#1).}\ \itshape
}{\par\vspace{0.3em}}
\title{Networking Protocols --- Detailed Lecture Notes}
\author{JKSSB Computer Awareness}
\date{\today}

\begin{document}
\maketitle

\section{Why protocols and addresses matter}
Computers on a network can understand each other only because they follow
agreed rules. A \textbf{protocol} is such a set of rules: it fixes what
messages look like and in which order the machines exchange them. Different
jobs need different protocols. Opening a web page, moving a file, sending an
email, translating a name into an address, and handing out addresses to new
devices are five different jobs, so they use five different protocols. The
exam tests each protocol by its \textbf{exact English name} together with the
single role it performs.

\begin{definitionbox}{Protocol}
A fixed set of rules governing communication between devices on a network,
covering message format and exchange order. TCP, IP, DNS, DHCP, HTTP, HTTPS,
FTP, Telnet, and SMTP are the protocols of this lesson.
\end{definitionbox}

\begin{center}
\begin{tabular}{p{0.24\textwidth}p{0.66\textwidth}}
\toprule
\textbf{Term} & \textbf{Meaning in this lesson} \\
\midrule
TCP & Transmission Control Protocol: reliable, ordered, connection-oriented delivery. \\
UDP & User Datagram Protocol: fast, connectionless, best-effort delivery. \\
IP & Internet Protocol: addressed, best-effort carriage of packets across networks. \\
ICMP & Internet Control Message Protocol: control and diagnostic messages only. \\
DNS & Domain Name System: translates names into IP addresses. \\
DHCP & Dynamic Host Configuration Protocol: automatically assigns addresses. \\
HTTP / HTTPS & HyperText Transfer Protocol (plain) and its Secure version for the web. \\
FTP & File Transfer Protocol: moves files between computers. \\
Telnet & Plain-text remote login protocol. \\
SMTP & Simple Mail Transfer Protocol: sends outgoing mail. \\
OSI model & The seven-layer reference model from Physical to Application. \\
Port & A number naming the service on a machine (HTTP 80, HTTPS 443). \\
\bottomrule
\end{tabular}
\end{center}

The rest of this lesson uses these terms in one consistent sense. \textbf{TCP}
and \textbf{UDP} are the two transport protocols. \textbf{IP} names the
addressing protocol with its two versions, IPv4 and IPv6. \textbf{ICMP} is the
helper that reports errors, never the carrier of user data. \textbf{DNS} and
\textbf{DHCP} are the two address services: one resolves names, the other
assigns addresses. \textbf{HTTP}, \textbf{HTTPS}, \textbf{FTP},
\textbf{Telnet}, and \textbf{SMTP} are application protocols, each with one
job and one well-known port.

\section{TCP and UDP}
The transport layer offers two opposite bargains. \textbf{TCP} (Transmission
Control Protocol) is connection-oriented: it sets up the connection, numbers
every segment, retransmits whatever is lost, and delivers the bytes complete
and in order. What TCP guarantees is \emph{ordering and correctness}; it does
not promise when the data will arrive. \textbf{UDP} (User Datagram Protocol)
is connectionless: it sends each datagram independently, with no setup, no
ordering, and no retransmission. What UDP offers in return is speed and low
overhead.

\begin{definitionbox}{TCP}
A connection-oriented transport protocol that delivers data reliably and in
order through sequencing, acknowledgement, and retransmission.
\end{definitionbox}

\begin{definitionbox}{UDP}
A connectionless transport protocol that sends independent datagrams quickly
with no delivery, ordering, or error-recovery guarantee.
\end{definitionbox}

Choose between them by what the task values most. A file, a web page, or an
email must arrive complete and correct, so it travels over TCP. Live voice, a
video stream, or a quick DNS lookup values speed over perfection, so it
travels over UDP, where a lost fragment is simply skipped.

\begin{example}
An item states, ``TCP guarantees fast delivery of every packet.'' This is
wrong. TCP guarantees that the bytes arrive complete and in order (PYQ-18);
speed is exactly what TCP sacrifices. The fast, best-effort protocol is UDP.
\end{example}

\section{IPv4 and IPv6}
The \textbf{Internet Protocol (IP)} carries addressed packets from the sender
to the receiver across networks. Each device holds an IP address, and routers
read that address to forward the packet hop by hop. IP by itself is
unreliable: delivery is best-effort, and it is TCP above it that adds
reliability.

\textbf{IPv4} uses 32-bit addresses written as four decimal octets, for
example 192.168.1.10. Its header is variable in length (20 to 60 bytes).
\textbf{IPv6} uses 128-bit addresses written as eight hexadecimal groups, for
example 2001:0db8::1, and its header is fixed at 40 bytes (PYQ-18, TRM-10).
The longer address space is the whole point of IPv6: it supplies vastly more
addresses than IPv4.

\begin{definitionbox}{IPv4}
The 32-bit version of the Internet Protocol, with dotted-decimal addresses
and a variable-length header.
\end{definitionbox}

\begin{definitionbox}{IPv6}
The 128-bit version of the Internet Protocol, with hexadecimal addresses and
a fixed 40-byte header.
\end{definitionbox}

\section{ICMP: the control channel}
\textbf{ICMP} (Internet Control Message Protocol) carries control and
diagnostic messages about the network itself: error reports such as
``destination unreachable'' and reachability probes. Its two familiar tools
are \textbf{ping}, which asks whether a host is reachable, and
\textbf{traceroute}, which shows the path packets take. ICMP transfers no
application data and uses no port numbers (TRAP-19). That last point is the
exam's favourite trap: any option giving ICMP a port, a file transfer, or a
mail delivery is wrong.

\section{DNS and DHCP}
\textbf{DNS} (Domain Name System) translates a human name such as
\texttt{jkssb.nic.in} into the IP address the network needs. Without it every
user would have to memorise numeric addresses. DNS queries travel mainly over
UDP on port 53, falling back to TCP for large zone transfers.

\textbf{DHCP} (Dynamic Host Configuration Protocol) automatically assigns an
IP address, subnet mask, default gateway, and DNS server address to each
device that joins the network. Without it every address would have to be
typed in by hand. The home router and the office network each run a DHCP
server; the joining laptop or phone acts as the DHCP client, using UDP ports
67 and 68.

\begin{definitionbox}{DNS}
The distributed naming service that resolves a domain name into its IP
address. It answers ``what number goes with this name?''
\end{definitionbox}

\begin{definitionbox}{DHCP}
The service that automatically assigns IP configuration (address, mask,
gateway, DNS server) to devices joining a network. It answers ``what address
may this newcomer use?''
\end{definitionbox}

\begin{example}
An item asks which service gives a newly connected office desktop its IP
address. Trace the pair: DNS resolves names the user already has, while DHCP
hands out addresses to newcomers. The answer is DHCP. If the stem instead
asks what turns \texttt{jkssb.nic.in} into an address, the answer is DNS.
\end{example}

\section{Application protocols: web, files, remote login, mail}
\textbf{HTTP} (HyperText Transfer Protocol, port 80) fetches web pages in
plain text. \textbf{HTTPS} (HyperText Transfer Protocol Secure, port 443) is
the same protocol with an SSL/TLS encryption layer, shown by
\texttt{https://} and the padlock. The S stands for Secure (WO 2026 Q76
territory).

\textbf{FTP} (File Transfer Protocol, port 21) moves files between computers;
it is neither email nor web browsing. \textbf{Telnet} (port 23) opens a
plain-text remote login session and is today replaced by the encrypted SSH.
\textbf{SMTP} (Simple Mail Transfer Protocol, port 25) sends outgoing mail to
the server. Receiving is a separate job handled by \textbf{POP3} (Post Office
Protocol version 3, port 110), which downloads mail to the local machine
(PYQ-16, TRAP-17), and \textbf{IMAP} (Internet Message Access Protocol, port
143), which keeps mail synced on the server.

\section{The seven OSI layers}
The \textbf{OSI model} (Open Systems Interconnection) is a seven-layer
reference frame for describing network communication:

\begin{center}
\begin{tabular}{p{0.10\textwidth}p{0.22\textwidth}p{0.58\textwidth}}
\toprule
\textbf{No.} & \textbf{Layer} & \textbf{Job} \\
\midrule
7 & Application & Services the user sees: HTTP, FTP, SMTP, DNS, DHCP. \\
6 & Presentation & Translation, encryption, compression. \\
5 & Session & Opens, manages, and closes sessions between machines. \\
4 & Transport & End-to-end delivery: TCP and UDP. \\
3 & Network & Addressing and routing: IP (IPv4, IPv6) and ICMP. \\
2 & Data Link & Frames on a single link: switch, MAC address. \\
1 & Physical & Bits on the wire, cable, or radio signal. \\
\bottomrule
\end{tabular}
\end{center}

Data moves down the layers on sending and up the layers on receiving, each
layer adding (or reading) its own header. For this exam, three placements
carry almost all the marks: the application names (HTTP, FTP, SMTP, DNS,
DHCP) at layer 7, TCP and UDP at layer 4, and IP and ICMP at layer 3.

\begin{example}
Opening a web page uses HTTP at the Application layer, TCP at the Transport
layer, and IP at the Network layer: one protocol per layer, working
together. An item asking ``which layer does IP belong to'' is answered by
the same table: the Network layer (layer 3).
\end{example}

\section{Port awareness}
An \textbf{IP address} names the machine; a \textbf{port number} names the
service on that machine. Ports 0--1023 are the well-known ports. The exam set
is: FTP 21, Telnet 23, SMTP 25, DNS 53, HTTP 80, POP3 110, IMAP 143, HTTPS
443. DHCP uses UDP ports 67 and 68, and ICMP uses no ports at all. The
recommended drill is to cover the service names and recall the numbers in the
chain 21--23--25--53--80--443.

\section{Exam retrieval and common confusions}
Seven distinctions prevent most errors on this topic.
\begin{enumerate}
  \item IPv6 is \textbf{128-bit}, never 256-bit; its header is fixed at 40
    bytes (TRAP-18, TRM-10).
  \item ICMP is \textbf{control/diagnostic} only: no application data, no
    ports (TRAP-19).
  \item TCP promises \textbf{order}, not timing; UDP promises neither
    (PYQ-18).
  \item \textbf{SMTP sends} mail; \textbf{POP3/IMAP receive} it (TRAP-17).
  \item \textbf{DNS resolves} names to addresses; \textbf{DHCP assigns}
    addresses to joining devices.
  \item HTTP is port \textbf{80} and plain; HTTPS is port \textbf{443} with
    SSL/TLS.
  \item TCP and UDP live at the \textbf{Transport} layer; IP and ICMP at the
    \textbf{Network} layer; HTTP, FTP, SMTP, DNS, and DHCP at the
    \textbf{Application} layer.
\end{enumerate}

\subsection*{Exercises}
\begin{enumerate}[label=\arabic*.]
  \item Distinguish TCP from UDP on reliability, connection setup, and one
    typical use each.
  \item Contrast IPv4 with IPv6 on address length, notation, and header size.
  \item Explain why ICMP is described as a control protocol and name its two
    common tools.
  \item Distinguish DNS from DHCP with one sentence on each role.
  \item List the seven OSI layers in order with the protocols of this lesson
    placed at their layers.
\end{enumerate}

\subsection*{Solutions}
\begingroup\small
\begin{enumerate}[label=\arabic*.]
  \item TCP is connection-oriented, reliable, and ordered: it numbers
    segments and retransmits losses; typical uses are web pages, email, and
    file transfer. UDP is connectionless and best-effort: no setup, ordering,
    or retransmission; typical uses are streaming, voice calls, and DNS
    queries.
  \item IPv4 is 32-bit, dotted-decimal (four octets), with a variable 20--60
    byte header. IPv6 is 128-bit, hexadecimal (eight groups), with a fixed
    40-byte header.
  \item ICMP carries control and diagnostic messages about the network
    (errors, reachability) rather than user data, and it uses no ports. Its
    tools are ping (reachability) and traceroute (path).
  \item DNS translates a domain name into its IP address. DHCP automatically
    assigns an IP address and related configuration to a device joining the
    network.
  \item 7 Application (HTTP, HTTPS, FTP, Telnet, SMTP, POP3, IMAP, DNS,
    DHCP); 6 Presentation; 5 Session; 4 Transport (TCP, UDP); 3 Network (IP,
    ICMP); 2 Data Link; 1 Physical.
\end{enumerate}
\endgroup
\end{document}
